\section{The IVAD Method}
\label{sec:method}

\subsection{Overview}

\method is designed to convert an agent trajectory into a selective,
evidence-constrained
diagnostic decision. Given a failed execution, the method (1) projects the raw
trace into a sanitized diagnostic context and addressable evidence catalog,
(2) generates a structured diagnosis, (3) independently re-parses its evidence
through deterministic and semantic verification channels, and (4) chooses
among acceptance, one bounded acquisition action, and abstention.
Figure~\ref{fig:ivad-system} shows how SpanVouch realizes projection,
structured diagnosis, dual-channel verification, bounded revision, human
decision, and manifest-bound persistence. The group-selective controller and
Level~1/2 acquisition actions are protocol components whose empirical claims
require qualifying calibration and provider artifacts; the engineering results
reported here do not substitute for those outcomes.

\subsection{Claim--Evidence Contract}

The claim--evidence contract prevents natural-language explanations from being
treated as self-authenticating evidence. Each causal claim names one or more
evidence identifiers, and every identifier resolves through a stable selector
to an immutable trajectory value and digest. Deterministic validation checks
the span tree, selector existence, canonical value, hash, duplicate references,
critical-span grounding, scope, and domain invariants. These checks are hard
constraints: a semantic score cannot override them.

Contract v1 has six roots:\newline
{\footnotesize\code{spanvouch.trace/1.0}; \code{spanvouch.diagnosis/1.0};\newline
\code{spanvouch.diagnostic-context/1.0};\newline
\code{spanvouch.verification/1.0}; \code{spanvouch.review/1.0};\newline
\code{spanvouch.artifact-manifest/1.0}.} Package versions and contract versions
are independent: a package release need not change a contract, and a contract
revision is not implied by a package release. Selectors, canonical values, and
hashes are therefore machine-checkable bindings. They do not establish that a
cited item is semantically relevant or that the cited items are causally
sufficient for a claim.

\subsection{Controlled Failure Separation}

The semantic verifier receives the structured claim, sanitized trajectory, and
allowed evidence catalog, but not the diagnoser's hidden reasoning or evaluator
labels. It judges whether the cited facts support the claimed failure and
responsibility, whether decisive counter-evidence is unaddressed, and whether
an alternative causal account remains viable. The experimental design varies
model identity, provider, prompt, retrieved summaries, and visible rationale
to measure which forms of separation reduce correlated false acceptance.
The B0--B5 design compares these conditions across paired LangGraph and
AutoGen scenario-template clusters. The offline matrix reported in this Technical Report
validates those interfaces with declared fake providers; estimating incremental
defect detection or correlated false acceptance requires real-provider outcomes.

\subsection{Group-Selective Risk Protocol}

Verification produces hard eligibility constraints and soft features such as
channel disagreement, evidence coverage, semantic support, and distributional
indicators. IVAD freezes a scalar error score, its direction, deterministic
tie-breaking, a finite candidate family $\Lambda$, target $\alpha$, confidence
level $1-\delta$, and minimum accepted-group count $m_{\min}$ before
calibration. The independently sampled unit is a preregistered task group, not
an individual report. A group incurs loss when any accepted report within it
has a frozen critical defect.

Calibration records $(K_\lambda^G,M_\lambda^G)$ for every candidate and applies
simultaneous one-sided exact-binomial bounds. The protocol selects the feasible
candidate with greatest calibration group acceptance, applies a frozen
tie-breaker, and evaluates the result once on untouched test data. If no
candidate meets both the risk bound and accepted-group minimum, IVAD reports no
feasible operating point. Section~\ref{sec:problem} states the full conditional
finite-sample guarantee and its assumptions. SpanVouch supplies the contracts,
artifact identities, and evaluation pipeline needed to execute this protocol;
the engineering results in this Technical Report do not claim empirical target-risk
attainment.

\subsection{One-Shot Layered Acquisition}

An acquisition action is triggered only by a typed evidence gap. The policy
selects the least costly permitted level, records the action and budget, allows
at most one diagnosis revision, and then repeats verification. The no-
acquisition and post-acquisition paths use a jointly frozen calibration
protocol over routing, action, revision, re-verification, and final acceptance.
If calibration labels select an acquisition policy, $\Lambda$ is the complete
frozen policy-by-threshold family covered by the simultaneous correction;
otherwise the policy is frozen on development data before calibration. No
pre-acquisition threshold is reused after acquisition. This design makes the
acceptance recovered by acquisition measurable against tokens, tool calls,
latency, and monetary cost.
